% Options for packages loaded elsewhere
\PassOptionsToPackage{unicode}{hyperref}
\PassOptionsToPackage{hyphens}{url}
\documentclass[
  ignorenonframetext,
]{beamer}
\newif\ifbibliography
\usepackage{pgfpages}
\setbeamertemplate{caption}[numbered]
\setbeamertemplate{caption label separator}{: }
\setbeamercolor{caption name}{fg=normal text.fg}
\beamertemplatenavigationsymbolsempty
% remove section numbering
\setbeamertemplate{part page}{
  \centering
  \begin{beamercolorbox}[sep=16pt,center]{part title}
    \usebeamerfont{part title}\insertpart\par
  \end{beamercolorbox}
}
\setbeamertemplate{section page}{
  \centering
  \begin{beamercolorbox}[sep=12pt,center]{section title}
    \usebeamerfont{section title}\insertsection\par
  \end{beamercolorbox}
}
\setbeamertemplate{subsection page}{
  \centering
  \begin{beamercolorbox}[sep=8pt,center]{subsection title}
    \usebeamerfont{subsection title}\insertsubsection\par
  \end{beamercolorbox}
}
% Prevent slide breaks in the middle of a paragraph
\widowpenalties 1 10000
\raggedbottom
\AtBeginPart{
  \frame{\partpage}
}
\AtBeginSection{
  \ifbibliography
  \else
    \frame{\sectionpage}
  \fi
}
\AtBeginSubsection{
  \frame{\subsectionpage}
}
\usepackage{iftex}
\ifPDFTeX
  \usepackage[T1]{fontenc}
  \usepackage[utf8]{inputenc}
  \usepackage{textcomp} % provide euro and other symbols
\else % if luatex or xetex
  \usepackage{unicode-math} % this also loads fontspec
  \defaultfontfeatures{Scale=MatchLowercase}
  \defaultfontfeatures[\rmfamily]{Ligatures=TeX,Scale=1}
\fi
\usepackage{lmodern}
\ifPDFTeX\else
  % xetex/luatex font selection
\fi
% Use upquote if available, for straight quotes in verbatim environments
\IfFileExists{upquote.sty}{\usepackage{upquote}}{}
\IfFileExists{microtype.sty}{% use microtype if available
  \usepackage[]{microtype}
  \UseMicrotypeSet[protrusion]{basicmath} % disable protrusion for tt fonts
}{}
\makeatletter
\@ifundefined{KOMAClassName}{% if non-KOMA class
  \IfFileExists{parskip.sty}{%
    \usepackage{parskip}
  }{% else
    \setlength{\parindent}{0pt}
    \setlength{\parskip}{6pt plus 2pt minus 1pt}}
}{% if KOMA class
  \KOMAoptions{parskip=half}}
\makeatother
\usepackage{color}
\usepackage{fancyvrb}
\newcommand{\VerbBar}{|}
\newcommand{\VERB}{\Verb[commandchars=\\\{\}]}
\DefineVerbatimEnvironment{Highlighting}{Verbatim}{commandchars=\\\{\}}
% Add ',fontsize=\small' for more characters per line
\newenvironment{Shaded}{}{}
\newcommand{\AlertTok}[1]{\textcolor[rgb]{1.00,0.00,0.00}{\textbf{#1}}}
\newcommand{\AnnotationTok}[1]{\textcolor[rgb]{0.38,0.63,0.69}{\textbf{\textit{#1}}}}
\newcommand{\AttributeTok}[1]{\textcolor[rgb]{0.49,0.56,0.16}{#1}}
\newcommand{\BaseNTok}[1]{\textcolor[rgb]{0.25,0.63,0.44}{#1}}
\newcommand{\BuiltInTok}[1]{\textcolor[rgb]{0.00,0.50,0.00}{#1}}
\newcommand{\CharTok}[1]{\textcolor[rgb]{0.25,0.44,0.63}{#1}}
\newcommand{\CommentTok}[1]{\textcolor[rgb]{0.38,0.63,0.69}{\textit{#1}}}
\newcommand{\CommentVarTok}[1]{\textcolor[rgb]{0.38,0.63,0.69}{\textbf{\textit{#1}}}}
\newcommand{\ConstantTok}[1]{\textcolor[rgb]{0.53,0.00,0.00}{#1}}
\newcommand{\ControlFlowTok}[1]{\textcolor[rgb]{0.00,0.44,0.13}{\textbf{#1}}}
\newcommand{\DataTypeTok}[1]{\textcolor[rgb]{0.56,0.13,0.00}{#1}}
\newcommand{\DecValTok}[1]{\textcolor[rgb]{0.25,0.63,0.44}{#1}}
\newcommand{\DocumentationTok}[1]{\textcolor[rgb]{0.73,0.13,0.13}{\textit{#1}}}
\newcommand{\ErrorTok}[1]{\textcolor[rgb]{1.00,0.00,0.00}{\textbf{#1}}}
\newcommand{\ExtensionTok}[1]{#1}
\newcommand{\FloatTok}[1]{\textcolor[rgb]{0.25,0.63,0.44}{#1}}
\newcommand{\FunctionTok}[1]{\textcolor[rgb]{0.02,0.16,0.49}{#1}}
\newcommand{\ImportTok}[1]{\textcolor[rgb]{0.00,0.50,0.00}{\textbf{#1}}}
\newcommand{\InformationTok}[1]{\textcolor[rgb]{0.38,0.63,0.69}{\textbf{\textit{#1}}}}
\newcommand{\KeywordTok}[1]{\textcolor[rgb]{0.00,0.44,0.13}{\textbf{#1}}}
\newcommand{\NormalTok}[1]{#1}
\newcommand{\OperatorTok}[1]{\textcolor[rgb]{0.40,0.40,0.40}{#1}}
\newcommand{\OtherTok}[1]{\textcolor[rgb]{0.00,0.44,0.13}{#1}}
\newcommand{\PreprocessorTok}[1]{\textcolor[rgb]{0.74,0.48,0.00}{#1}}
\newcommand{\RegionMarkerTok}[1]{#1}
\newcommand{\SpecialCharTok}[1]{\textcolor[rgb]{0.25,0.44,0.63}{#1}}
\newcommand{\SpecialStringTok}[1]{\textcolor[rgb]{0.73,0.40,0.53}{#1}}
\newcommand{\StringTok}[1]{\textcolor[rgb]{0.25,0.44,0.63}{#1}}
\newcommand{\VariableTok}[1]{\textcolor[rgb]{0.10,0.09,0.49}{#1}}
\newcommand{\VerbatimStringTok}[1]{\textcolor[rgb]{0.25,0.44,0.63}{#1}}
\newcommand{\WarningTok}[1]{\textcolor[rgb]{0.38,0.63,0.69}{\textbf{\textit{#1}}}}
\setlength{\emergencystretch}{3em} % prevent overfull lines
\providecommand{\tightlist}{%
  \setlength{\itemsep}{0pt}\setlength{\parskip}{0pt}}
% Auto-generated by infrastructure/rendering/_pdf_latex_helpers.py::extract_math_font_preamble.
% Minimal header passed to Pandoc via -H so Beamer slides pick up the same math font
% as the combined-PDF path. Adjust the manuscript's preamble.md to change the font globally.
\usepackage{unicode-math}
\setmathfont{latinmodern-math.otf}
\usepackage{bookmark}
\IfFileExists{xurl.sty}{\usepackage{xurl}}{} % add URL line breaks if available
\urlstyle{same}
\hypersetup{
  hidelinks,
  pdfcreator={LaTeX via pandoc}}

\author{\texorpdfstring{}{}}
\date{}

\begin{document}

\section{Manuscript Syntax Reference}\label{manuscript-syntax-reference}

\begin{frame}[fragile,allowframebreaks]{Manuscript Syntax Reference}
Formatting and syntax conventions for manuscript Markdown files in the
\texttt{code\_project} exemplar.
\end{frame}

\begin{frame}[fragile,allowframebreaks]{Citation Syntax (Pandoc)}
\protect\phantomsection\label{citation-syntax-pandoc}
\begin{Shaded}
\begin{Highlighting}[]
\CommentTok{\textless{}!{-}{-} Single citation {-}{-}\textgreater{}}
\CommentTok{[}\OtherTok{@knuth1997}\CommentTok{]}

\CommentTok{\textless{}!{-}{-} Multiple citations {-}{-}\textgreater{}}
\CommentTok{[}\OtherTok{@knuth1997; @cormen2009}\CommentTok{]}

\CommentTok{\textless{}!{-}{-} Citation with locator {-}{-}\textgreater{}}
\CommentTok{[}\OtherTok{@knuth1997, pp. 42{-}45}\CommentTok{]}

\CommentTok{\textless{}!{-}{-} Narrative citation {-}{-}\textgreater{}}
\NormalTok{@knuth1997 showed that...}
\end{Highlighting}
\end{Shaded}

All citation keys must exist in \texttt{references.bib}. The pipeline
will flag undefined keys as errors.
\end{frame}

\begin{frame}[fragile,allowframebreaks]{Equation Environments}
\protect\phantomsection\label{equation-environments}
\begin{Shaded}
\begin{Highlighting}[]
\CommentTok{\textless{}!{-}{-} Numbered equation with label {-}{-}\textgreater{}}
\NormalTok{\textbackslash{}begin\{equation\}}
\NormalTok{\textbackslash{}label\{eq:gradient\}}
\NormalTok{\textbackslash{}nabla f(x) = Ax {-} b}
\NormalTok{\textbackslash{}end\{equation\}}

\CommentTok{\textless{}!{-}{-} Reference in text {-}{-}\textgreater{}}
\NormalTok{As shown in Equation \textbackslash{}ref\{eq:gradient\}, the gradient...}
\end{Highlighting}
\end{Shaded}

Pandoc-crossref resolves \texttt{\{\#eq:label\}} and \texttt{@eq:label}
during rendering. Do \textbf{not} use raw
\texttt{\textbackslash{}label\{\}} / \texttt{\textbackslash{}ref\{\}}
--- the Markdown filter handles this.
\end{frame}

\begin{frame}[fragile,allowframebreaks]{Figure References}
\protect\phantomsection\label{figure-references}
\begin{Shaded}
\begin{Highlighting}[]
\CommentTok{\textless{}!{-}{-} Figure with label {-}{-}\textgreater{}}
\AlertTok{![Convergence plot showing objective value vs iteration](figures/convergence\_plot.png)}\NormalTok{\{\#fig:convergence width=80\%\}}

\CommentTok{\textless{}!{-}{-} Reference in text {-}{-}\textgreater{}}
\NormalTok{@fig:convergence demonstrates...}
\end{Highlighting}
\end{Shaded}

\begin{itemize}
\tightlist
\item
  Images must exist in \texttt{output/figures/} at render time
\item
  Use \texttt{width=} or \texttt{height=} to control sizing (prevents
  float-too-large warnings)
\item
  Captions should be descriptive (they appear in the PDF)
\end{itemize}
\end{frame}

\begin{frame}[fragile,allowframebreaks]{Table References}
\protect\phantomsection\label{table-references}
\begin{Shaded}
\begin{Highlighting}[]
\CommentTok{\textless{}!{-}{-} Table with label {-}{-}\textgreater{}}
\PreprocessorTok{|}\NormalTok{ Step Size }\PreprocessorTok{|}\NormalTok{ Iterations }\PreprocessorTok{|}\NormalTok{ Converged }\PreprocessorTok{|}
\PreprocessorTok{|{-}{-}{-}{-}{-}{-}{-}{-}{-}{-}{-}|{-}{-}{-}{-}{-}{-}{-}{-}{-}{-}{-}|{-}{-}{-}{-}{-}{-}{-}{-}{-}{-}{-}|}
\PreprocessorTok{|}\NormalTok{ 0.01      }\PreprocessorTok{|}\NormalTok{ 412       }\PreprocessorTok{|}\NormalTok{ Yes       }\PreprocessorTok{|}

\NormalTok{: Performance comparison across step sizes \{\#tbl:performance\}}

\CommentTok{\textless{}!{-}{-} Reference in text {-}{-}\textgreater{}}
\NormalTok{@tbl:performance shows...}
\end{Highlighting}
\end{Shaded}
\end{frame}

\begin{frame}[fragile,allowframebreaks]{Preamble Injection}
\protect\phantomsection\label{preamble-injection}
\texttt{preamble.md} contains raw LaTeX injected before the document
body:

\begin{Shaded}
\begin{Highlighting}[]
\CommentTok{{-}{-}{-}}
\AnnotationTok{header{-}includes:}
\CommentTok{  {-} \textbackslash{}usepackage\{amsmath\}}
\CommentTok{  {-} \textbackslash{}tolerance=800}
\CommentTok{  {-} \textbackslash{}hfuzz=2pt}
\CommentTok{{-}{-}{-}}
\end{Highlighting}
\end{Shaded}

\begin{itemize}
\tightlist
\item
  Preamble is parsed by
  \texttt{infrastructure.rendering.latex\_utils.py}
\item
  Changes to tolerance/hfuzz affect line-breaking globally
\item
  Do \textbf{not} duplicate package imports already in the
  infrastructure renderer
\end{itemize}
\end{frame}

\begin{frame}[fragile,allowframebreaks]{BibTeX Entry Format}
\protect\phantomsection\label{bibtex-entry-format}
\begin{Shaded}
\begin{Highlighting}[]
\VariableTok{@article}\NormalTok{\{}\OtherTok{knuth1997}\NormalTok{,}
  \DataTypeTok{author}\NormalTok{  = \{Donald E. Knuth\},}
  \DataTypeTok{title}\NormalTok{   = \{The Art of Computer Programming\},}
  \DataTypeTok{journal}\NormalTok{ = \{Addison{-}Wesley\},}
  \DataTypeTok{year}\NormalTok{    = \{1997\},}
  \DataTypeTok{volume}\NormalTok{  = \{1\}}
\NormalTok{\}}
\end{Highlighting}
\end{Shaded}

\begin{itemize}
\tightlist
\item
  Keys must be lowercase alphanumeric with optional underscores
\item
  Author names use \texttt{\{Last,\ First\}} or \texttt{\{First\ Last\}}
  format
\item
  All entries must have at minimum: \texttt{author}, \texttt{title},
  \texttt{year}
\end{itemize}
\end{frame}

\begin{frame}[fragile,allowframebreaks]{Section Numbering}
\protect\phantomsection\label{section-numbering}
\begin{Shaded}
\begin{Highlighting}[]
\NormalTok{00\_abstract.md    → Section 0 (unnumbered in PDF)}
\NormalTok{01\_introduction.md → Section 1}
\NormalTok{02\_methodology.md  → Section 2}
\NormalTok{03\_results.md      → Section 3}
\NormalTok{04\_conclusion.md   → Section 4}
\end{Highlighting}
\end{Shaded}

Files are assembled in lexicographic order by
\texttt{infrastructure.rendering.pdf\_renderer.py}. The \texttt{00\_}
prefix ensures the abstract renders first.
\end{frame}

\begin{frame}[fragile,allowframebreaks]{Prose Conventions}
\protect\phantomsection\label{prose-conventions}
\begin{itemize}
\tightlist
\item
  No ``In summary'' or ``In conclusion'' at section ends (RASP standard)
\item
  Use active voice for methodology descriptions
\item
  Use explicit file paths when referencing code:
  \texttt{src/optimizer.py}, not ``the optimizer module''
\item
  Keep paragraphs focused --- one idea per paragraph
\end{itemize}
\end{frame}

\begin{frame}[allowframebreaks]{See Also}
\protect\phantomsection\label{see-also}
\begin{itemize}
\tightlist
\item
  \url{AGENTS.md} --- RASP protocol and AI agent constraints
\item
  \href{../docs/rendering_pipeline.md}{rendering\_pipeline.md} --- Full
  rendering flow
\item
  \url{preamble.md} --- Active LaTeX preamble
\end{itemize}
\end{frame}

\end{document}
